\section{Marco teórico}

\subsection{Habilidades blandas: fundamentos teóricos}

\begin{itemize}
    \item \textbf{Competencias:} De acuerdo con el Ministerio de Educación Nacional, es el “conjunto de conocimientos, actitudes, disposiciones y habilidades (cognitivas, socioafectivas y comunicativas), relacionadas entre sí para facilitar el desempeño flexible y con sentido de una actividad en contextos relativamente nuevos y retadores. Por lo tanto, la competencia implica conocer, ser y saber hacer” \cite{Ministerio2024}.
    
    \item \textbf{Habilidades blandas:} Las habilidades blandas son competencias no técnicas que facilitan un desempeño eficaz en contextos sociales y laborales diversos. A diferencia de las habilidades técnicas, las habilidades blandas no dependen de la interacción con máquinas y son aplicables en una variedad de contextos profesionales, implican capacidades interpersonales e intrapersonales. Se consideran componentes subyacentes de la competencia de un individuo y están relacionadas con su motivación, rasgos de personalidad y conocimiento \cite{Marin-Zapata2022}.
\end{itemize}

\subsection{Teorías del aprendizaje}
Las teorías del aprendizaje son útiles para comprender cómo se adquieren y se modifican conocimientos, creencias, conductas, actitudes, habilidades y destrezas. 
\\
\\
A continuación, se presentan las teorías del aprendizaje seleccionadas por su relevancia y aplicabilidad al diseño de los módulos implementados en este proyecto.

\subsubsection{Aprendizaje autorregulado}
El aprendizaje autorregulado (Self-Regulated Learning, SRL) es una perspectiva teórica que explica cómo los estudiantes gestionan de manera activa y deliberada sus propios procesos cognitivos, motivacionales y conductales para alcanzar metas de aprendizaje. Este modelo ha sido desarrollado principalmente por Barry J. Zimmerman, quien lo define como “el proceso autodirigdo mediante el cual los estudiantes transforman sus habilidades mentales en habilidades académicas” \cite{Zimmerman2002}. Esta perspectiva señala que aprender no consiste únicamente en poseer ciertas capacidades cognitivas, si no en saber cómo organizar, monitorear, regular y evalur el propio desempeño durante una tarea académica.
\\
\\
Zimmerman señala que los aprendices autorregulados se caracterizan por adoptar un rol activo en su formación: identifican sus fortalezas y limitaciones, establecen metas claras, seleccionan estrategias aadecuadas, supervisan su ejecución y ajustan su comportamiento en función de los resultados obtenidos. De acuerdo con el autor “el aprendizaje es una actividad que los estudiantes realizan por sí mismos de manera proactiva, y no un evento pasivo que ocurre como reacción a la enseñanza” \cite{Zimmerman2002}.
\\
\\
El modelo de aprendizaje autorregulado se organiza en tres fases cíclicas, las cuales interactúan entre sí y permiten que el estudiante mejore progresivamente su capacidad de aprender de manera autónoma:

\paragraph{1. Fase de premeditación (forethought)}

Esta fase ocurre antes de iniciar una tarea e involucra procesos de análisis y preparación:

\begin{itemize}
    \item \textbf{Análisis de la tarea:} Incluye el establecimiento de metas específicas y la planificación estratégica. Zimmerman destaca que los estudiantes que fijan metas cercanas y bien definidas suelen mostrar mayores niveles de autoeficacia y persistencia.
    
    \item \textbf{Creencias motivacionales:} Comprende la autoeficacia, las expectativas de resultado, el interés intrínseco y la orientación hacia el aprendizaje. Estas creencias influyen directamente en la disposición con la que el estudiante enfrenta la actividad.
\end{itemize}

\paragraph{2. Fase de desempeño (performance)}

Corresponde a los procesos que ocurren durante la ejecución de la tarea:

\begin{itemize}
    \item \textbf{Autocontrol:} Implica el uso de estrategias cognitivas, metacognitivas y conductuales, tales como la autoinstrucción, la organización del entorno de estudio y la regulación de la atención.
    
    \item \textbf{Autoobservación:} Se refiere al monitoreo continuo del propio desempeño, el registro de avances y la identificación de errores. Zimmerman sostiene que incluso prácticas simples de auto-monitoreo pueden generar mejoras espontáneas en el rendimiento académico.
\end{itemize}

\paragraph{3. Fase de autorreflexión (self-reflection)}

Incluye los procesos posteriores a la ejecución de la tarea:

\begin{itemize}
    \item \textbf{Juicio personal:} Consiste en comparar los resultados obtenidos con las metas previamente establecidas, así como en analizar las causas del éxito o fracaso.
    
    \item \textbf{Reacciones adaptativas o defensivas:} Implica ajustar estrategias, modificar creencias o tomar decisiones para futuras tareas. Atribuciones causales centradas en factores controlables, como el uso de estrategias, fortalecen la motivación; en contraste, atribuciones basadas en la falta de capacidad suelen conducir a la desmotivación.
\end{itemize}

Zimmerman plantea que estas fases se articulan como un ciclo continuo, ya que la reflexión y evaluación de una experiencia influyen directamente en la planificación de la siguiente. De este modo, el aprendizaje autorregulado se concibe como un proceso dinámico y progresivo, que puede mejorar gradualmente a través de la enseñanza y la práctica continua.

\subsubsection{Metacognición}
La metacognición es el conjunto de conocimientos y experiencias que las personas poseen sobre sus propios procesos cognitivos, así como la capacidad de supervisarlos, regularlos y controlarlos durante actividades intelectuales. Este concepto fue introducido formalmente por John H. Flavell en su obra \textbf{“Metacognition and Cognitive Monitoring”} \cite{Flavell1979}, donde indica que estos procesos resultan esenciales en tareas como la comprensión, la memoria, la comunicación, la resolución de problemas y distintas formas de autocontrol cognitivo.
\\
\\
Flavell explica que el monitoreo cognitivo está constituido de cuatro elementos centrales \cite{Flavell1979}: El conocimiento metacognitivo, las experiencias metacognitivas, las metas o tareas y las acciones o estrategias. Estos componentes operan de manera articulada para permitir que las personas regulen de forma efectiva su pensamiento y aprendizaje.

\paragraph{1. Conocimiento metacognitivo}

Es el conjunto de creencias y saberes que una persona posee acerca de sí misma y de los demás como aprendices, de las características de las tareas cognitivas y de las estrategias que pueden utilizarse para afrontarlas. Flavell organiza este conocimiento en tres categorías:

\begin{itemize}
    \item \textbf{Persona:} incluye creencias sobre uno mismo y sobre otros como agentes cognitivos, considerando diferencias individuales y propiedades universales del funcionamiento mental.
    \item \textbf{Tarea:} comprende el entendimiento de las características de una actividad (dificultad, familiaridad, estructura) y de cómo estas influyen en el desempeño.
    \item \textbf{Estrategia:} alude al conocimiento de los procedimientos cognitivos o metacognitivos disponibles, así como su eficacia en distintos contextos.
\end{itemize}

Flavell señala que este conocimiento puede activarse de forma consciente o automática y no siempre es exacto, lo cual afecta la calidad del monitoreo cognitivo.

\paragraph{2. Experiencias metacognitivas:}
Son los estados cognitivos y afectivos que surgen antes, durante o después de una tarea intelectual. Pueden manifestarse como dudas, sensaciones de comprensión o falta de ella, percepciones de progreso o dificultad, o sentimientos de seguridad respecto a la ejecución.
\\
\\
Estas experiencias cumplen tres funciones principales:

\begin{itemize}
    \item \textbf{Ajustar metas o tareas:} por ejemplo, volver a leer un texto al percibir falta de comprensión.
    \item \textbf{Modificar el conocimiento metacognitivo:} actualizar creencias o corregir percepciones previas.
    \item \textbf{Activar estrategias cognitivas o metacognitivas:} como autoevaluar la comprensión o repetir información difícil.
\end{itemize}

\paragraph{3. Metas o tareas}

Se refieren a los objetivos que guían la actividad cognitiva. Según Flavell, la naturaleza de la meta (su claridad, dificultad o relevancia) determina qué estrategias son necesarias y qué grado de monitoreo se debe aplicar para completarla con éxito.

\paragraph{4. Acciones o estrategias}

Son los procedimientos utilizados para alcanzar metas cognitivas o metacognitivas. Flavell distingue entre:

\begin{itemize}
    \item \textbf{Estrategias cognitivas:} orientadas a mejorar el desempeño directo en una tarea (por ejemplo, releer para memorizar).
    \item \textbf{Estrategias metacognitivas:} orientadas a evaluar o regular el propio desempeño (por ejemplo, autoformular preguntas para comprobar comprensión).
\end{itemize}


\subsubsection{Teoría del constructivismo}
La teoría del constructivismo, atribuida a Jean Piaget, sostiene que el aprendizaje es un proceso activo en el que los individuos construyen su propio conocimiento a través de la interacción con el entorno y la reinterpretación de sus experiencias. De tal forma, se reconoce que cada persona es responsable de elaborar su comprensión del mundo a partir de sus experiencias previas, conocimientos previos y percepciones individuales.
\\
\\
El conocimiento es una construcción del ser humano: cada persona percibe la realidad, la organiza y le da sentido en forma de constructos, lo que contribuye a la edificación de un todo coherente que da sentido y unicidad a la realidad \cite{OrtizGranja2015}. Por lo tanto, al ser el individuo un participante activo en su propio proceso de aprendizaje es crucial fomentar la autonomía, recordando la importancia de aprender por cuenta propia.

\subsubsection{Teoría del aprendizaje social}
La teoría del aprendizaje social de Albert Bandura refuerza la idea de que el aprendizaje es el resultado de la observación al entorno social y la imitación. En su libro Social Learning Theory (1977) menciona: “Afortunadamente, en la mayoría de los humanos el comportamiento se aprende de manera observacional a través del modelado: al observar a otros se forma una idea de cómo se realizan los nuevos comportamientos, y en ocasiones posteriores, esta información codificada sirve como guía para la acción” \cite{Bandura1977}.
\\
\\
En su investigación, Bandura describe el aprendizaje social en cuatro principios: El primero indica que para aprender, es necesario prestar atención y estar enfocado; el segundo señala que la información recién adquirida debe ser internalizada y almacenada como recuerdo; el tercero implica la reproducción de la información retenida; y finalmente, el cuarto principio destaca la importancia de tener motivación, el cual es originado al observar las consecuencias, ya sean recompensas o castigos, que otros experimentan por sus acciones o decisiones.

\subsubsection{Teoría del aprendizaje experiencial}
La teoría del aprendizaje experiencial de David Kolb se cimenta en desarrollos y trabajos teóricos previos de autores como Jean Piaget, John Dewey y Kurt Lewin. Esta teoría establece la importancia de la experiencia en el proceso de aprendizaje. El desarrollo de David Kolb explora los procesos cognitivos asociados al abordaje y procesamiento de las experiencias y describe los diferentes estilos individuales de aprendizaje \cite{GomezPawelek2011}.
\\
\\
Kolb, propone seguir un proceso de cuatro etapas denominado “Ciclo de Kolb” para que exista un aprendizaje efectivo. Este ciclo empieza con la experiencia concreta, que se adquiere al hacer algo, luego se lleva a cabo la reflexión sobre los resultados obtenidos (etapa de observación reflexiva), después se extraen conclusiones o generalizaciones (etapa de conceptualización abstracta) a partir de la experiencia y por último, se pone en práctica las conclusiones formuladas para orientar nuestras acciones en situaciones futuras (etapa de experimentación activa).

\subsubsection{Teoría del aprendizaje significativo}
De acuerdo con el psicólogo educativo David Ausubel, el aprendizaje significativo se da cuando una nueva información es integrada y comprendida en función de un conocimiento anterior, en lugar de ser memorizada de manera aislada. Se produce así una interacción entre los nuevos contenidos y los elementos relevantes ya presentes en la estructura cognitiva. La presencia de ideas, conceptos o proposiciones inclusivas, claras y disponibles en la mente del aprendiz es lo que dota de significado al nuevo aprendizaje \cite{PalmeroRodriguez2024}.

\subsubsection{Teoría del aprendizaje por descubrimiento}
Jerome S. Bruner, plantea la importancia del descubrimiento personal en el aprendizaje, para que sea significativo. Según Bruner, el individuo adquiere conocimiento cuando es descubierto por sí mismo o por su propio discernimiento, ya que involucra la exploración activa, la experimentación y la resolución de problemas por parte de éste. Así mismo, con este tipo de aprendizaje es más probable que la información sea conservada.
\\
\\
Para que se produzca un aprendizaje por descubrimiento, es indispensable que se planteen objetivos específicos que guíen la investigación del individuo.  Estos objetivos actúan como puntos de referencia que orientan el proceso de descubrimiento y facilitan la consecución de resultados concretos. Además, es fundamental que exista una familiarización con las herramientas utilizadas para el descubrimiento de información, ya que estas proporcionan los medios necesarios para llevar a cabo cualquier investigación de manera efectiva. 
\\
\\
Por último, Bruner sugiere que el tanto investigar como experimentar propicia la fijación de conocimientos y es la manera más natural de adquirir competencias. También ayuda a la motivación del aprendiz, a fomentar el uso de la creatividad en la resolución de problemas, a estimular la intuición e impulsa el aprender a aprender. 

\subsection{Procesamiento del Lenguaje Natural (PLN)}
El Procesamiento del Lenguaje Natural (PLN) es una rama de la inteligencia artificial que estudia cómo los sistemas computacionales pueden representar, comprender, analizar y generar lenguaje humano, tanto en su forma escrita como hablada. Según lo planteado por Jurafky y Martin en “Speech and Language Processing” \cite{JurafskyMartin2023}, el PLN integra principios de la lingüística, modelos matemáticos, técnicas de aprendizaje automático e ingeniería de sistemas para procesar grandes volúmenes de información lingüística, modelar su estructura y significado, así como interpretar su uso en distintos contextos comunicativos.

\subsubsection{Modelos de lenguaje de gran escala (Large Language Models, LLM)}

Los modelos de lenguaje de gran escala (Large Language Models, LLM) constituyen una rama emergente dentro del procesamiento del lenguaje natural, caracterizada por arquitecturas profundas basadas en transformadores que contienen decenas de miles de millones de parámetros y que han sido entrenadas en vastos corpus textuales. Según el recurso de Elastic, un LLM es “un modelo entrenado con algoritmos de aprendizaje profundo capaz de realizar una amplia gama de tareas de procesamiento de lenguaje natural” \cite{ElasticLLM}.
\\
\\
El funcionamiento básico de estos modelos implica tres etapas fundamentales:
\begin{itemize}
  \item \textbf{Pre-entrenamiento:} el modelo es entrenado sin supervisión en grandes cantidades de datos lingüísticos, lo que le permite inferir patrones estadísticos del lenguaje \cite{ElasticLLM}.
  \item \textbf{Ajuste fino (fine-tuning):} para tareas específicas, el modelo se adapta con datos etiquetados o ejemplos, mejorando su rendimiento en dominios particulares \cite{ElasticLLM}.
  \item \textbf{Prompt-tuning o aprendizaje en contexto (few-shot/zero-shot):} el modelo es guiado mediante instrucciones o ejemplos mínimos para realizar nuevas tareas sin entrenamiento adicional extenso \cite{ElasticLLM}.
\end{itemize}

Además, los LLM suelen emplear tres componentes técnicos principales: capas de embedding para representar palabras o tokens en vectores, mecanismos de atención para detectar dependencias de largo alcance y capas feed-forward que procesan la información codificada. \cite{ElasticLLM}.
\\
\\
Estos modelos han sido aplicados con éxito en múltiples ámbitos: generación de texto, chatbots conversacionales, generación de código, análisis de sentimientos, motores de búsqueda mejorados y sistemas de recomendación. No obstante, también enfrentan desafíos significativos como las alucinaciones, los sesgos, los riesgos de privacidad y los altos requerimientos computacionales \cite{ElasticLLM}.
\\
\\
Dentro del contexto de este proyecto, los LLM constituyen un fundamento tecnológico para la generación automática de retroalimentación, la interpretación del lenguaje de los usuarios, la producción de contenido educativo personalizado y el soporte de agentes conversacionales inteligentes que intervienen en los módulos diseñados.


\subsection{Teoría general de sistemas}

La teoría general de los sistemas (TGS), propuesta por Ludwig Von Bertalanffy, ampliamente aplicada en diversos campos de conocimiento, es definida como el estudio interdisciplinario de los sistemas y sus interrelaciones dentro de sistemas más complejos. Una premisa fundamental de esta teoría es: “el todo es más que la suma de las partes” (Aristóteles), es decir, al examinar de manera holística la interacción entre diferentes componentes de un sistema, pueden surgir características de éste que no son evidentes si se observarán cada una de sus partes de forma independiente.
\\
\\
La TGS se distingue por su perspectiva integradora, donde se considera importante la interacción y los conjuntos que a partir de ella se generan, ofreciendo un ambiente adecuado para la interrelación y comunicación entre especialistas. Con ello, contradice la mirada fragmentaria y lineal del enfoque reduccionista, método con el cual los fenómenos se estudian de una forma independiente como si fueran una realidad dividida \cite{GutierrezGomez2013}.

\subsection{Principios de diseño y usabilidad}

La usabilidad es un aspecto central en el diseño de sistemas interactivos y se refiere al grado en que un producto permite a los usuarios alcanzar objetivos específicos con eficacia, eficiencia y satisfacción. Nielsen, establece que la usabilidad no es un rasgo único o unidimensional, sino un concepto compuesto por cinco dimensiones que permiten evaluar de forma integral la experiencia del usuario: \textit{aprendibilidad}, \textit{eficiencia}, \textit{memorabilidad}, \textit{errores} y \textit{satisfacción}. Estas dimensiones describen, respectivamente, qué tan fácil resulta para un usuario aprender a usar un sistema, qué tan rápido puede ejecutar tareas, cómo recuerda su funcionamiento tras un periodo de no uso, qué tan propenso es a cometer errores y qué tan satisfecho se siente con la experiencia.
\\
\\
\textbf{Heurísticas de usabilidad de Nielsen.}  
Nielsen propone diez heurísticas que funcionan como principios generales para el diseño de interfaces centradas en el usuario \cite{NielsenMolich1990}. Estas heurísticas son ampliamente utilizadas en ingeniería de usabilidad y diseño de experiencias digitales:

\begin{enumerate}
    \item \textbf{Visibilidad del estado del sistema.} 
    El sistema debe mantener informados a los usuarios sobre lo que está ocurriendo, proporcionando retroalimentación oportuna y comprensible.

    \item \textbf{Correspondencia entre el sistema y el mundo real.}
    La interfaz debe emplear conceptos, terminología y convenciones familiares para los usuarios, siguiendo el orden lógico y natural propio del mundo real.

    \item \textbf{Control y libertad del usuario.}
    El sistema debe permitir que los usuarios deshagan y rehagan acciones fácilmente, así como navegar o salir de estados no deseados sin dificultad.

    \item \textbf{Consistencia y estándares.}
    La interfaz debe mantener coherencia en su diseño y seguir convenciones establecidas, evitando que el usuario tenga que adivinar si diferentes elementos realizan acciones similares.

    \item \textbf{Prevención de errores.}
    El diseño debe anticipar y minimizar la aparición de errores, ya sea restringiendo entradas no válidas o proporcionando mecanismos que reduzcan la posibilidad de equivocación.

    \item \textbf{Reconocimiento antes que recordación.}
    El sistema debe reducir la carga de memoria del usuario al mostrar opciones, elementos y acciones relevantes, evitando que dependa de recordar información entre distintas partes de la interfaz.

    \item \textbf{Flexibilidad y eficiencia de uso.}
    La interfaz debe ser usable tanto para usuarios novatos como expertos, incorporando atajos, personalización u otros mecanismos que permitan acelerar tareas frecuentes.

    \item \textbf{Diseño estético y minimalista.}
    El contenido presentado debe ser claro y esencial, evitando información irrelevante o excesiva que pueda distraer o sobrecargar cognitivamente al usuario.

    \item \textbf{Ayuda a los usuarios a reconocer, diagnosticar y recuperarse de errores.}
    Los mensajes de error deben expresarse en lenguaje claro, indicando el problema, su causa y posibles soluciones concretas para recuperarse.

    \item \textbf{Ayuda y documentación.}
    Aunque lo ideal es que el sistema sea intuitivo, debe ofrecer documentación accesible, fácil de buscar y enfocada en describir pasos concretos para completar tareas.
\end{enumerate}

Estas heurísticas constituyen una guía práctica para evaluar y mejorar interfaces, y son aplicables tanto en sistemas web como en aplicaciones móviles, plataformas educativas o herramientas basadas en inteligencia artificial.
\\
\\
\textbf{Principios de diseño centrado en el usuario:}  
Norman, en \textit{The Design of Everyday Things} \cite{Norman2013}, destaca que el diseño centrado en el usuario se fundamenta en la comprensión profunda de las necesidades, objetivos y limitaciones de los usuarios. Dentro de este enfoque se resaltan cuatro principios clave:

\begin{itemize}
    \item \textbf{Comprensión de los usuarios:} análisis de necesidades, contexto, limitaciones y comportamientos.
    \item \textbf{Diseño iterativo:} ciclos de prototipado, evaluación y refinamiento progresivo.
    \item \textbf{Evaluación continua:} pruebas con usuarios durante todas las fases del desarrollo.
    \item \textbf{Diseño inclusivo:} consideración de diversidad de usuarios, capacidades y contextos.
\end{itemize}

El enfoque centrado en el usuario es especialmente relevante para sistemas educativos y herramientas basadas en inteligencia artificial, donde la facilidad de uso y la claridad de las interacciones influyen directamente en la experiencia y el aprendizaje del usuario.
